% !TEX root = ../matan.tex

\section{Теорема о почленном интегрировании функционального ряда.}

\begin{Theorem}{Почленное интегрирование равномерно сходящегося ряда}{}
	Пусть $I = [a, b]$, функции $f_k \colon I \to \RR$ интегрируемы по Риману на $I$, и ряд $\sum_{k=1}^{\infty} f_k$ сходится равномерно на $I$ к функции $f$. Тогда $f$ интегрируема по Риману на $I$ и
	\[
	\int_a^b f(x)\,dx = \sum_{k=1}^{\infty} \int_a^b f_k(x)\,dx.
	\]
\end{Theorem}

\begin{proof}
	Обозначим частичные суммы и остатки
	\[
	S_n(x) = \sum_{k=1}^{n} f_k(x), \qquad r_n(x) = f(x) - S_n(x) = \sum_{k=n+1}^{\infty} f_k(x).
	\]
	Из равномерной сходимости ряда следует $\|r_n\|_\infty = \sup_{x \in I} |r_n(x)| \to 0$.

	\textbf{Интегрируемость $f$.} Для функции $h$ и отрезка $J \subset I$ обозначим её колебание через $\omega(h; J) = \sup_J h - \inf_J h$. По критерию Римана функция интегрируема тогда и только тогда, когда для любого $\eps > 0$ существует разбиение $\tau = \{a = x_0 < x_1 < \cdots < x_N = b\}$ с
	\[
	\sum_{m=0}^{N-1} \omega(h; [x_m, x_{m+1}])\,(x_{m+1} - x_m) < \eps.
	\]
	Зафиксируем $\eps > 0$ и выберем $n_0$ так, что $\|r_{n_0}\|_\infty < \dfrac{\eps}{4(b - a)}$. Сумма $S_{n_0}$ конечна и состоит из интегрируемых функций, значит интегрируема; поэтому найдётся разбиение $\tau$ с
	\[
	\sum_{m=0}^{N-1} \omega(S_{n_0}; [x_m, x_{m+1}])\,(x_{m+1} - x_m) < \frac{\eps}{2}.
	\]
	На каждом отрезке $f = S_{n_0} + r_{n_0}$, а колебание субаддитивно, поэтому
	\[
	\omega(f; [x_m, x_{m+1}]) \le \omega(S_{n_0}; [x_m, x_{m+1}]) + \omega(r_{n_0}; [x_m, x_{m+1}]) \le \omega(S_{n_0}; [x_m, x_{m+1}]) + 2\|r_{n_0}\|_\infty.
	\]
	Суммируя по отрезкам разбиения,
	\[
	\sum_{m=0}^{N-1} \omega(f; [x_m, x_{m+1}])\,(x_{m+1} - x_m) < \frac{\eps}{2} + 2\|r_{n_0}\|_\infty (b - a) < \frac{\eps}{2} + \frac{\eps}{2} = \eps.
	\]
	Значит, $f$ интегрируема по Риману на $I$.

	\textbf{Формула для интеграла.} Так как $f = S_n + r_n$,
	\[
	\int_a^b f(x)\,dx = \int_a^b S_n(x)\,dx + \int_a^b r_n(x)\,dx,
	\]
	причём
	\[
	\Bigl| \int_a^b r_n(x)\,dx \Bigr| \le (b - a)\,\|r_n\|_\infty \to 0.
	\]
	Поэтому $\int_a^b f = \lim_{n \to \infty} \int_a^b S_n$. Но в силу линейности интеграла
	\[
	\int_a^b S_n(x)\,dx = \int_a^b \sum_{k=1}^{n} f_k(x)\,dx = \sum_{k=1}^{n} \int_a^b f_k(x)\,dx,
	\]
	и, переходя к пределу при $n \to \infty$, получаем
	\[
	\int_a^b f(x)\,dx = \sum_{k=1}^{\infty} \int_a^b f_k(x)\,dx. \qedhere
	\]
\end{proof}

\begin{Remark}{}{}
	Равномерность сходимости существенна: при только поточечной сходимости предел и интеграл, вообще говоря, переставлять нельзя (площадь под графиками может не стремиться к площади под предельной функцией).
\end{Remark}
